\\maketitle
\\begin{abstract}
    本文针对2022年CUMCM A题化学反应动力学建模问题，建立了基于常微分方程组的反应动力学模型。首先根据质量作用定律推导反应速率方程；其次采用最小二乘法估计反应速率常数；最后通过灵敏度分析和数值模拟验证模型可靠性。
    
    \\textbf{针对问题一}，我们建立了反应动力学微分方程组。设反应物浓度为$c_i(t)$，则：
    $$\\frac{dc_i}{dt} = \\sum_j \\nu_{ij} k_j \\prod_l c_l^{\\alpha_{jl}}$$
    其中$\\nu_{ij}$为化学计量系数，$k_j$为速率常数。
    
    \\textbf{针对问题二}，我们设计了参数估计算法。采用Levenberg-Marquardt方法最小化目标函数：
    $$\\min \\sum_i \\sum_t (c_i^{model}(t) - c_i^{exp}(t))^2$$
    
    \\textbf{针对问题三}，我们进行了模型验证与灵敏度分析。通过残差分析和交叉验证，验证了模型的预测精度。局部灵敏度分析表明，速率常数$k_1$和$k_3$对产物浓度影响最大。
    
    本研究为化学反应过程优化控制提供了可靠的数学模型。
    
    \\keywords{反应动力学\quad 常微分方程\quad 参数估计\quad Levenberg-Marquardt\quad 灵敏度分析}
\\end{abstract}
